\section{Introduction}
LLM agents now plan, call tools, write code, and operate in software repositories with increasing autonomy. This progress has produced benchmarks for software engineering and cybersecurity agents, including bug fixing, web exploitation, CTF solving, vulnerability reproduction, patching, and exploit construction~\cite{swebench,programbench,cybergym,cvebench,secbench,bountybench,exploitgym,exploitbench,krepro}. A common assumption underlies many of these tasks: the world has already been prepared for the agent. The agent receives a repository, build script, patch, container, harness, running service, or executable target.

Known-vulnerability reproduction in the wild often begins in a less structured state. A security analyst may start with only a CVE identifier and must discover the affected product, locate the correct artifact, recover a runnable environment, and validate vulnerability behavior against the real target. IoT firmware vulnerabilities make this gap especially sharp. Firmware is often distributed as a proprietary binary blob, compiled for MIPS or ARM, packed in vendor-specific formats, dependent on missing NVRAM state and device files, and reachable only after the vulnerable service is rehosted under emulation. Thus firmware reproduction is not merely exploit writing; it is an end-to-end artifact reconstruction and validation task.

We introduce \sys, a benchmark that evaluates whether LLM agents can move from a bare CVE identifier toward auditable evidence of IoT firmware vulnerability reproduction. This setting follows the spirit of ProgramBench~\cite{programbench}, which moves beyond localized coding tasks to evaluate full-program reconstruction. \sys{} similarly moves beyond prepared vulnerability environments to evaluate the full artifact-to-execution pipeline. The benchmark stresses capabilities that prepared-environment benchmarks largely hide: open-world information gathering, artifact provenance checking, cross-architecture binary analysis, long-horizon tool orchestration, recovery from tool failures, and grounded validation.

The key design choice in \sys{} is that it does not assume a prebuilt executable oracle for each vulnerability. Instead, each task is grounded in a manually inspected \emph{public-evidence dossier}: a reference packet of public vulnerability facts and expected evidence types. Agents receive only the CVE identifier, not the dossier. After execution, their plans, traces, workspaces, scripts, logs, reports, and network or process evidence are scored against the dossier. This design reflects the realistic starting point of from-scratch reproduction while avoiding over-claiming that every IoT CVE has been independently re-executed by the benchmark authors.

\sys{} also addresses a methodological problem: plausible-looking reproductions are easy to fabricate. When rehosting is difficult, agents often produce substitutes such as a Python web server that mimics a router endpoint, a synthetic C program that crashes, or a static report with a copied PoC. Such outputs can look successful under weak evaluation, but they do not exercise the real firmware. \sys{} therefore uses non-binary, evidence-grounded scoring. It scores both the reproduction plan and the executed task across six stages, while enforcing real-target gating for service rehosting and vulnerability triggering.

Our experimental design contains two linked components. First, agents attempt 30 CVEs with 5 models and 3 trials, producing 450 full traces and workspaces. Second, a rubric-constrained scoring skill evaluates all 450 runs; post-hoc human analysis audits the scoring outputs, suspicious high scores, all R5/R6 claims, and the final failure-mode taxonomy. The central question is not whether an agent writes a convincing final report, but whether its artifacts support each claimed stage of reproduction.

Our contributions are:
\begin{itemize}
\item We formulate CVE-only IoT firmware vulnerability reproduction as a from-scratch, long-horizon agent benchmark.
\item We construct a manually curated 30-CVE public-evidence dataset balanced across buffer overflow, command injection, and authentication bypass tasks.
\item We propose a non-binary, evidence-grounded scoring framework for planning and execution across six dependent stages, with real-target gating for rehosting and triggering.
\item We design a 450-run evaluation and post-hoc human audit pipeline to characterize where agents lose auditable evidence and when they substitute simulations for real firmware execution.
\end{itemize}
